%% ============================================================
%% Blockchain in Healthcare Today (BHTY)
%% Opinions / Perspectives / Point of View
%% Compile with: pdflatex manuscript_draft_v2.tex  (x2)
%% ============================================================
\documentclass[12pt,a4paper]{article}

%% -- Encoding & fonts -----------------------------------------
\usepackage[T1]{fontenc}
\usepackage[utf8]{inputenc}
\usepackage{mathptmx}
\usepackage{microtype}

%% -- Page geometry -------------------------------------------
\usepackage[top=2.54cm, bottom=2.54cm,
            left=2.54cm,  right=2.54cm]{geometry}

%% -- Line spacing --------------------------------------------
\usepackage{setspace}
\onehalfspacing

%% -- Paragraph style -----------------------------------------
\usepackage{parskip}
\setlength{\parskip}{6pt}
\setlength{\parindent}{0pt}

%% -- Section heading style -----------------------------------
\usepackage{titlesec}
\titleformat{\section}[block]
  {\normalfont\bfseries\fontsize{13}{15}\selectfont}{}{0pt}{}
\titlespacing*{\section}{0pt}{14pt}{6pt}
\titleformat{\subsection}[block]
  {\normalfont\bfseries\itshape\fontsize{12}{14}\selectfont}{}{0pt}{}
\titlespacing*{\subsection}{0pt}{10pt}{4pt}

%% -- Lists ---------------------------------------------------
\usepackage{enumitem}
\setlist[enumerate]{nosep, left=0pt}
\setlist[itemize]{nosep, left=0pt}

%% -- Hyperlinks & URLs ---------------------------------------
\usepackage[colorlinks=true, linkcolor=black,
            urlcolor=blue, citecolor=black]{hyperref}
\usepackage{url}
\urlstyle{same}

%% -- Colour (article-type label) ----------------------------
\usepackage{xcolor}
\definecolor{bhtyblue}{RGB}{0,84,166}

%% -- Vancouver-style superscript citations -------------------
\usepackage[numbers,super,sort&compress]{natbib}
\renewcommand{\bibsection}{}

%% -- Page header/footer -------------------------------------
\usepackage{fancyhdr}
\setlength{\headheight}{14pt}
\pagestyle{fancy}
\fancyhf{}
\fancyhead[L]{\small\textit{Blockchain in Healthcare Today}}
\fancyhead[R]{\small\thepage}
\renewcommand{\headrulewidth}{0.4pt}

\date{}

%% ============================================================
\begin{document}
\thispagestyle{empty}

%% ---- ARTICLE TYPE LABEL ------------------------------------
\begin{center}
  {\fontsize{10}{12}\selectfont\bfseries\color{bhtyblue}
   OPINIONS / PERSPECTIVES / POINT OF VIEW}
\end{center}

\vspace{6pt}

%% ---- TITLE -------------------------------------------------
\begin{center}
  {\fontsize{16}{20}\selectfont\bfseries
   Governing Trustworthy Agentic Healthcare AI: A
   Calibration-First Framework with Distributed Ledger
   Accountability}
\end{center}

\vspace{10pt}

%% ---- AUTHORS & AFFILIATIONS --------------------------------
\begin{center}
  {\fontsize{12}{14}\selectfont
  Michael O. Eniolade\textsuperscript{1}%
  \,\href{https://orcid.org/0009-0002-0859-3050}{\texttt{0009-0002-0859-3050}}%
  }
\end{center}

\begin{center}
  {\fontsize{10}{12}\selectfont
  \textsuperscript{1}University of the Cumberlands,
  Williamsburg, KY, USA}
\end{center}

\vspace{4pt}
\begin{center}
  {\fontsize{10}{12}\selectfont
  \textbf{Corresponding Author:} Michael O.~Eniolade \quad
  Email: \href{mailto:meniolade20593@ucumberlands.edu}%
  {meniolade20593@ucumberlands.edu}}
\end{center}

\vspace{4pt}
\begin{center}
  {\fontsize{10}{12}\selectfont
  \textbf{DOI:} [Assigned upon acceptance]}
\end{center}

\vspace{8pt}

%% ---- KEYWORDS ----------------------------------------------
\noindent{\fontsize{10}{12}\selectfont
\textbf{Keywords:} agentic AI; healthcare governance;
probability calibration; blockchain; clinical decision support}

\vspace{10pt}
\hrule height 0.6pt
\vspace{10pt}

%% ============================================================
%% ABSTRACT
%% ============================================================
\section*{Abstract}

Healthcare AI has entered a new phase. Autonomous systems capable
of reasoning across multiple steps, querying external tools, and
initiating clinical actions are already operating in hospitals and
health systems around the world. These agentic systems raise
governance questions of a different order than those addressed by
existing frameworks. TRIPOD+AI, DECIDE-AI, and the FDA guidance
on AI/ML-based software as a medical device were all designed
around a simpler picture: a model makes one prediction, a
clinician reviews it. Pipeline-level reliability, compounding
errors across sequential decision steps, and the accountability
gap left when autonomous systems act faster than any human review
cycle are simply outside their scope.

This perspective proposes a five-layer trustworthiness governance
framework for agentic healthcare AI. The five layers address:
calibration integrity and uncertainty quantification at each agent
decision point; subgroup fairness auditing across demographic
groups; verifiable model provenance; patient-controlled consent
management for autonomous AI actions; and scored deployment
readiness with ongoing monitoring. Each layer carries a concrete,
measurable threshold: Expected Calibration Error at or below
0.10, conformal prediction coverage at or above 0.90, and
subgroup ECE gap no wider than 0.05. These are numbers, not
aspirations.

Blockchain and distributed ledger technologies (DLT) are not a
peripheral concern here. They are the accountability
infrastructure that makes this governance framework enforceable
rather than advisory. Immutable audit trails, smart
contract-encoded constraints, and decentralized consent
registries close the gap that traditional logging leaves open:
the inability to prove, after the fact, what an AI agent did and
why. The framework connects to active empirical research on
clinical AI calibration, conformal prediction, and deployment
readiness, and each layer is designed to be tested in subsequent
work.

\vspace{8pt}
\hrule height 0.6pt
\vspace{10pt}

%% ============================================================
%% PLAIN LANGUAGE SUMMARY
%% ============================================================
\section*{Plain Language Summary}

\textbf{What this paper is about.}
AI systems in healthcare are becoming agentic. They make
sequences of decisions on their own, without waiting for a human
to approve each step. An AI agent might monitor a patient's vital
signs, detect a deterioration pattern, and automatically alert
the care team or escalate the case. These systems are already in
clinical use, and their capabilities are growing.

\textbf{The problem.}
The rules and guidelines we use to make sure healthcare AI is
safe were built for simpler tools that make one decision at a
time. Those older rules do not address what happens when AI makes
many decisions in a row, how errors build on each other across
steps, or who is responsible when an autonomous system causes
harm through a chain of automated actions.

\textbf{What we propose.}
This paper proposes five requirements for any agentic AI system
deployed in healthcare: (1)~its probability predictions should be
accurate and well-calibrated at every decision step; (2)~it
should work equitably across different groups of patients;
(3)~its training history and data sources should be traceable and
verifiable; (4)~patients should be able to specify what AI is and
is not permitted to do on their behalf; and (5)~it should pass a
structured readiness assessment before going live, with ongoing
monitoring after deployment.

\textbf{Why blockchain helps.}
Blockchain technology creates permanent, tamper-proof digital
records. We argue it provides the technical foundation to make
these five requirements enforceable and independently verifiable,
rather than just written-down guidelines. Blockchain-based audit
trails, smart contracts, and decentralized consent registries
turn governance rules into operational constraints built into the
system itself, not into a policy binder on a shelf.

\textbf{Who this is for.}
Clinical informaticists, health system leaders, AI developers,
and policymakers who want a practical, technically grounded
governance framework for deploying agentic AI responsibly in
healthcare settings.

\vspace{8pt}
\hrule height 0.6pt
\vspace{10pt}

%% ============================================================
%% SECTION 1 -- INTRODUCTION
%% ============================================================
\section*{Introduction}

\subsection*{The Agentic Shift in Healthcare AI}

Healthcare AI used to be relatively straightforward to evaluate.
Train a model on historical data, validate it on a held-out test
set, report discrimination metrics, and check whether predictions
clear a performance threshold. One patient, one input, one
prediction. Clinical decision-support systems working this way
function like sophisticated calculators: a clinician asks a
question, the model returns a number.

That picture is changing. Agentic AI systems are autonomous,
goal-directed software entities capable of multi-step reasoning,
tool use, planning, and sequential action. They do not stop after
one response. They pursue goals across multiple steps, decompose
tasks, invoke external databases, retrieve records, make
intermediate decisions, and produce outputs that trigger further
automated actions downstream. Clinical deployments already include systems that autonomously
monitor deterioration signals across multiple data streams and
initiate escalation protocols. Diagnostic reasoning agents
iteratively query laboratory and imaging results to refine a
differential diagnosis. Medication safety agents cross-check
orders against renal function values, interaction databases, and
allergy records. Care coordination systems adjust clinical
workflows without a human initiating each
step.\cite{rajpurkar2022} A 2025 systematic review of AI agents
in clinical medicine documented twenty studies spanning discrete
clinical tasks including diagnostic reasoning, medication
management, and evidence retrieval, with agent systems
outperforming their base model counterparts by a median of
53~percentage points in single-agent tool-calling studies,
underscoring both the clinical promise and the governance urgency
of agentic deployment.\cite{gorenshtein2025}

\subsection*{Why This Changes the Governance Problem}

The governance implications of this shift are structural, not
incremental. When AI works as a single-inference tool,
accountability is relatively clear: a clinician requests a
prediction, evaluates it, and decides whether to act. The AI
functions as a consultant whose opinion is accepted or rejected.
When AI works as an agent, that structure breaks down.
The agent is no longer a consultant. It is a participant in care
delivery. Actions occur at speeds that preclude real-time human
review at each step. Errors compound: a wrong intermediate
decision at step two shapes everything that follows. Responsibility
spreads across the AI system, its developers, the deploying
institution, and the clinicians nominally overseeing the process.

Pinning down who is accountable for a poor outcome becomes
genuinely difficult. This is not a hypothetical concern. It is
the logical structure of autonomous pipeline deployment.

\subsection*{The Governance Gap}

Existing clinical AI governance frameworks were not designed for
this environment. TRIPOD+AI\cite{collins2024tripod} provides
reporting standards for model development and validation but
addresses models, not pipelines. DECIDE-AI\cite{vasey2022} guides
evaluation of clinical decision-support interventions but draws no
distinction between a single-inference tool and a multi-step
agent. The FDA's AI/ML-based Software as a Medical Device (SaMD)
framework\cite{fda2021} addresses pre-market validation and
post-market surveillance, but it was conceived before agentic
deployment became clinically common. The EU AI
Act\cite{euaiact2024} imposes requirements for high-risk AI
systems in healthcare but provides little operational guidance for
autonomous pipeline governance. None of these frameworks specifies
how to measure, record, or verify trustworthiness at the level of
individual agent decisions inside a running clinical system.

We propose a five-layer trustworthiness governance framework to
address that gap, and argue that distributed ledger technologies
provide the accountability infrastructure needed to make it
operational and verifiable at scale.

%% ============================================================
%% SECTION 2 -- WHAT MAKES AGENTIC AI DIFFERENT
%% ============================================================
\section*{What Makes Agentic Healthcare AI Different}

Good governance starts with an accurate picture of what is being
governed. We identify four properties of agentic AI systems that
create challenges absent from traditional clinical prediction
models.

\subsection*{Sequential Decision-Making and Error Compounding}

A traditional clinical AI model makes one prediction from one
input. An agentic system makes a chain of decisions where each
output feeds the next input. This introduces error compounding:
an overconfident probability estimate at step two, left
undetected, propagates into step four as though it were reliable.
Calibration matters at every link in the chain, not only at the
final output.

Current calibration evaluation methods were built for
single-output models. A model's Expected Calibration Error (ECE)
is computed across its test set predictions as a
whole.\cite{guo2017} In an agentic pipeline, each decision node
carries its own calibration profile. Aggregate pipeline
reliability reflects the product of individual node reliabilities.
A pipeline of five individually well-calibrated components still
produces systematically biased outputs if calibration error drifts
in the same direction at each step.\cite{angelopoulos2023}

\subsection*{Opacity of Reasoning Chains}

Single-inference models are amenable to well-established
interpretability tools. SHAP values decompose a prediction into
feature contributions. Calibration curves expose probability
reliability. Permutation importance identifies which inputs drive
a classification. Agentic reasoning chains, particularly those
involving large language model components, iterative retrieval,
or multi-agent orchestration, resist these approaches. The
intermediate states between initial input and final recommendation
are often partially or fully inaccessible, and their contribution
to the final output cannot always be recovered.

When an agentic system produces a harmful recommendation,
investigators face a genuine epistemic problem: which step
introduced the error? Was it a miscalibrated prediction, a biased
training distribution, or an inappropriate action at a pipeline
junction? Without a durable record of intermediate states,
post-incident analysis becomes speculation rather than
investigation.

\subsection*{Speed Beyond Human Review Cycles}

Agentic AI acts at speeds that outpace human review. A triage
agent processing 50 patient deterioration signals per hour cannot
be supervised by a clinician reviewing each action in real time.
Meaningful oversight in this environment requires something
different: governance embedded in the system's design through
pre-deployment constraints, hardcoded boundaries, and automatic
escalation triggers, then verified through post-hoc audit rather
than concurrent human review.\cite{choudhury2020}

This is a design requirement, not a workaround.

\subsection*{Diffuse Moral Responsibility}

When a single prediction model leads to patient harm, the
responsibility chain is at least recognizable: the model
developer, the deploying institution, the clinician who acted on
the output. Agentic pipelines dissolve that chain. Multiple AI
components, orchestration layers, external data sources, and
automated actions each contribute to an outcome, but no single
party controls the full sequence.\cite{zuiderwijk2021}

The result is not just a legal problem. Diffuse responsibility
creates practical barriers to learning from failures and to
deterring unsafe deployment in the first place.

%% ============================================================
%% SECTION 3 -- TRUSTWORTHINESS GAP
%% ============================================================
\section*{The Trustworthiness Gap in Current Governance
Frameworks}

\subsection*{What Existing Frameworks Address}

Several frameworks have emerged to govern clinical AI quality.
TRIPOD+AI\cite{collins2024tripod} extends the TRIPOD reporting
guideline to cover AI-specific transparency requirements: model
architecture, training data provenance, calibration reporting, and
external validation. DECIDE-AI\cite{vasey2022} guides evaluation
of AI-based clinical decision-support interventions, with a focus
on clinical utility alongside statistical performance. The FDA
SaMD framework\cite{fda2021} establishes a risk-based regulatory
pathway for AI-enabled medical devices, including pre-market
validation and post-market surveillance requirements. The EU AI
Act\cite{euaiact2024} classifies healthcare AI as high-risk and
imposes transparency, documentation, human oversight, and
robustness requirements.

Each of these frameworks represents genuine progress. None of
them was designed for pipeline-level governance of autonomous
systems.

\subsection*{What They Miss}

\textbf{Calibration is underspecified.}
TRIPOD+AI recommends calibration reporting but sets no threshold
standards. A 2024 systematic review and meta-analysis of
prediction models for incident chronic kidney disease derived or
validated in community-based electronic health records found
calibration reporting so deficient that pooled meta-analysis was
not possible. Risk of bias was high in 64\% of included models,
and no clinical utility analyses or impact studies were identified
for any model.\cite{haris2024} An agent decision node that
produces unreliable probability estimates passes existing
governance checks without any flag.

\textbf{Uncertainty goes unreported.}
Fewer than 4\% of clinical AI studies report any uncertainty
measure alongside their predictions.\cite{prinzi2025} For a
standalone model, this is a quality gap. For an agentic system
where uncertain predictions propagate into subsequent automated
actions, it is a safety gap. An agent unaware of its own
uncertainty at step three has no basis to pause and request human
review before acting.

\textbf{Fairness evaluation stops at the model level.}
Where fairness auditing appears in clinical AI governance, it
typically asks whether subgroup performance gaps exist in a
model's test set. It does not ask whether bias accrues across
pipeline steps, or whether demographic disparities in agent
recommendations emerge from the interaction of multiple
components that are individually unbiased.\cite{mittelstadt2016}

\textbf{Audit records are informal and mutable.}
No current governance framework specifies what a sufficient audit
record for an AI agent's actions looks like. Server logs, model
outputs, and database entries are institution-specific, editable,
and not independently verifiable. In a regulatory dispute or
adverse event review, the absence of immutable, standardized
records makes accountability analysis unreliable.

\textbf{Consent frameworks predate autonomous action.}
Existing informed consent mechanisms in healthcare AI address
patient consent to data use and model deployment at the
institutional level. They were not designed for a world where an
AI agent initiates actions on a patient's behalf, step by step,
across a clinical episode.\cite{mittelstadt2016} Meaningful
consent in that context requires something fundamentally different
from a signature at admission.

%% ============================================================
%% SECTION 4 -- FIVE-LAYER FRAMEWORK
%% ============================================================
\section*{A Five-Layer Trustworthiness Governance Framework}

We propose a five-layer framework built to address each gap
identified above. Every layer carries measurable criteria and
applies at two points: pre-deployment assessment and continuous
monitoring of running systems.

\subsection*{Layer 1: Reliability Evaluation}

Reliability evaluation asks whether the probability outputs from
each agent decision node are trustworthy enough to drive automated
downstream action.

\textbf{Calibration assessment} requires computing the Expected
Calibration Error (ECE) for each decision node using empirical
frequency analysis across probability bins.\cite{guo2017}
We propose an ECE of $\leq 0.10$ as a minimum deployment
threshold, consistent with emerging quality standards in clinical
AI. The Brier Score and Brier Skill Score serve as complementary
summary metrics; Guo et~al.\cite{guo2017} show how both expose
systematic overconfidence in deep learning models. Calibration
must be evaluated on the actual distribution of inputs the
deployed agent encounters, not only on the original training
distribution. Real clinical inputs are heterogeneous,
missing-data-rich, and often quite different from what was in the
training set.

\textbf{Uncertainty quantification} using split conformal
prediction\cite{angelopoulos2023} provides distribution-free
coverage guarantees at each decision node. Instead of a single
probability estimate, conformal prediction produces a prediction
set: a set of possible outputs guaranteed to contain the true
label with probability $1-\alpha$, where $\alpha$ is the chosen
miscoverage rate. We propose a coverage guarantee of $\geq 0.90$
for clinical agentic AI. When a prediction set contains multiple
labels, the agent should defer to human review rather than proceed
autonomously. Ambiguity is a signal to pause, not to guess.

\textbf{External validation} is required before any deployment.
Calibration metrics should be recomputed on a prospectively
collected or geographically distinct external dataset to quantify
calibration drift: the degradation in probability reliability
when a model meets a new clinical
environment.\cite{haris2024}

\subsection*{Layer 2: Fairness Auditing}

Fairness auditing asks whether agent reliability is equitably
distributed across patient subgroups.

Calibration should be stratified by age group (e.g., below 65
vs.\ 65 and older), sex, race and ethnicity where available, and
clinically relevant comorbidity groups such as diabetes status
and chronic kidney disease stage. We propose a maximum allowable
ECE gap of 0.05 between any two demographic subgroups as a
deployment gate. Where subgroup calibration gaps exceed that
threshold, targeted recalibration on underrepresented groups or
restricted deployment to populations with demonstrated
calibration should follow.

False negative disparity deserves particular attention. Missed
diagnoses fall disproportionately on groups already
disadvantaged in the health system.\cite{mittelstadt2016}
An agent that is well-calibrated on average but
systematically underestimates risk for elderly patients or
patients from racial minority groups deepens existing inequities
rather than closing them.

\subsection*{Layer 3: Provenance Verification}

Provenance verification asks whether the lineage of an agent's
models and training data is independently verifiable.

\textbf{Model versioning} requires each deployed agent decision
node to carry a verifiable identifier: a cryptographic hash of
the model weights and architecture, recorded at deployment. Any
change to the model, including recalibration, fine-tuning, or
architectural modification, generates a new identifier and a new
deployment record.

\textbf{Training data documentation} requires that dataset
characteristics be documented and verifiable: size, demographic
composition, data source, collection period, and preprocessing
steps. An independent reviewer should be able to assess from this
documentation whether the training distribution matches the
deployment environment.

\textbf{Deployment history} records when each model version was
deployed, to which clinical contexts, and under what oversight
conditions. When a patient is harmed, investigators need to know
exactly which model version was running at the time of the
incident, what it was trained on, and whether its deployment
conditions matched its validation conditions.

\subsection*{Layer 4: Consent Management}

Consent management governs the scope of autonomous AI action to
which a patient has agreed.

Current healthcare AI consent frameworks address consent to data
use and model deployment at the institutional level. Agentic AI
requires something more specific. A patient should be able to
specify whether they consent to autonomous AI actions that affect
their care plan, whether they require human review before
AI-initiated actions take effect, and whether they retain the
right to override an agent recommendation at any step.

Expressing, storing, and enforcing those preferences requires a
consent architecture capable of updating as the patient's care
context changes. A form signed at admission does not meet this
bar for a system where AI agents act across multiple care
episodes, data systems, and clinical teams.

\subsection*{Layer 5: Deployment Readiness and Ongoing
Monitoring}

Deployment readiness scoring brings together the preceding four
layers into a structured pre-deployment gate and a continuous
post-deployment monitoring protocol.

Before deployment, each agent is assessed against a checklist of
eight criteria: (1)~discrimination adequacy (AUROC $\geq 0.85$
on external validation); (2)~calibration adequacy
(ECE $\leq 0.10$ on external validation); (3)~calibration
stability (calibration drift $\leq 0.05$ across internal and
external datasets); (4)~uncertainty coverage (conformal coverage
$\geq 0.90$ on external validation); (5)~coverage stability
(coverage drift $\leq 0.05$ across datasets); (6)~prediction
interpretability (singleton rate $\geq 0.70$, meaning at least
70\% of instances receive an unambiguous prediction);
(7)~subgroup calibration equity (ECE gap across subgroups
$\leq 0.05$); and (8)~transparency and reproducibility (full
code, data pipeline, and model documentation publicly or
institutionally accessible).\cite{collins2024eval}

After deployment, sequential ECE monitoring with statistical
process control alerts tracks calibration drift over time. A
statistically significant ECE increase above the deployment
threshold triggers a mandatory review cycle. If the problem is
not resolved, the system is decommissioned.

%% ============================================================
%% SECTION 5 -- DLT AS ACCOUNTABILITY INFRASTRUCTURE
%% ============================================================
\section*{Distributed Ledger Technologies as Accountability
Infrastructure}

\subsection*{Why Traditional Logging Falls Short}

The five-layer framework generates a substantial evidence base:
calibration metrics, fairness audits, provenance records, consent
documents, deployment readiness scores, and ongoing monitoring
data. For that evidence to serve accountability in a regulatory
review, adverse event investigation, or litigation, it must be
durable, verifiable, and independent of any single institution's
control.

Traditional server logs and database records fail on all three
counts. They are mutable: administrators with access can alter
records. They are siloed: records from different components of a
multi-institutional agentic pipeline live in incompatible systems
with no common reference point. They are not independently
verifiable: an institution claiming its AI agent behaved correctly
at a specific time cannot prove its own internal records are
unmodified.

In healthcare environments where agents act across institutional
boundaries, a diagnostic agent querying a regional radiology
system, a triage agent integrating data from multiple hospitals,
no single institution's logging architecture provides an
authoritative audit record of the full pipeline.

\subsection*{Blockchain-Based Audit Trails}

Blockchain and distributed ledger technologies address these
limitations through properties well-suited to the accountability
requirements of agentic healthcare AI. As blockchain and AI
converge as complementary pillars of digital health
transformation, their integration for governance and
accountability has emerged as a priority challenge for healthcare
systems globally.\cite{dershem2025}

\textbf{Immutability} means a transaction written to the ledger
cannot be altered without network consensus, making post-hoc
modification of audit records computationally and
organizationally prohibitive. Each significant agent decision: a
risk score above a clinical threshold, an automated order, a care
plan modification, goes onto the chain as a tamper-resistant
record of what the agent did, when, on what inputs, and with what
output.

\textbf{Verifiability} means any authorized party independently
confirms a record is unmodified since it was written. A
regulatory auditor, a patient's legal representative, or a
clinical ethics committee verifies an agent's audit trail
without relying on the deploying institution's attestation.

\textbf{Decentralization} means the audit record belongs to no
single party. In a multi-institutional pipeline, a shared
distributed ledger lets each participant contribute to and verify
the record without any one party having unilateral control over
what the record says.\cite{ullah2024}

\textbf{Smart contracts} encode governance rules as executable
code running on the ledger. Rather than relying on institutional
policy, a smart contract enforces governance rules automatically:
if an agent's conformal prediction set contains multiple labels,
the contract blocks autonomous action and logs a human escalation
event on-chain. Ullah et~al.\ demonstrate this approach in the
EHR domain, implementing purpose-based access control policies
through smart contracts that establish comprehensive, immutable
audit trails recording transaction details, access purposes,
digital signatures, and timestamps, providing accountability for
all data access events while eliminating dependence on trusted
intermediaries.\cite{ullah2024} Governance rules become
operational constraints with independently verifiable
enforcement, not aspirational policies.

\subsection*{Decentralized Consent Management}

Patient consent for autonomous AI action is a natural fit for
blockchain-based management. A consent record written to a
distributed ledger is patient-controlled, institution-independent,
and auditable. A patient who consents to AI-assisted triage but
not AI-initiated medication changes encodes that preference in a
consent record queried by any pipeline component before it acts.

Smart contracts enforce consent boundaries in real time. When an
agent proposes an action, the contract checks whether that action
class falls within the patient's recorded consent scope and blocks
the action if not, generating a logged exception for human review.
Consent becomes a live constraint on system behavior, not a
document filed at admission.

\subsection*{Verifiable Model Provenance}

The provenance requirements in Layer~3 are directly implementable
through on-chain model registries. When a new model version goes
live, its cryptographic hash: computed from model weights,
architecture specification, and training data characteristics, is
written to the ledger. Any query to the agent can then verify in
real time that the model responding is the registered, validated
version and not a modified derivative.

This matters especially in environments where models are updated
periodically. A post-deployment recalibration that improves
performance on a new patient population sometimes degrades
performance on a previously studied subgroup. An on-chain
registry makes every version change traceable and independently
verifiable. Liu and Hu's systematic review of blockchain-AI
integration in healthcare, spanning 1,221 publications from 2015
to 2025, identifies model provenance, traceability, and smart
contract-enforced data governance as foundational requirements
for building a trustworthy and intelligent healthcare ecosystem,
with blockchain providing the immutability and auditability that
AI-generated clinical outputs require.\cite{liu2026}

\subsection*{Federated Learning for Privacy-Preserving Distributed
Governance}

Governance at scale requires access to diverse patient populations
for calibration validation and fairness auditing across multiple
hospitals, health systems, and jurisdictions. Federated learning
trains and calibrates models across distributed datasets without
centralizing sensitive patient data. Distributed ledger
technologies provide the coordination, verification, and audit
infrastructure for the federated process.\cite{rieke2020}

Each participating institution trains on local data, contributes
model statistics to the federation, and receives a globally
calibrated model whose performance has been verified across the
full population diversity of the federation. The ledger records
each institution's participation, the aggregation protocol, and
the resulting model provenance, producing accountability for the
entire federated governance process.

%% ============================================================
%% SECTION 6 -- IMPLEMENTATION CHALLENGES
%% ============================================================
\section*{Implementation Challenges}

The framework is architecturally coherent, but implementation
requires confronting several real obstacles.

\subsection*{Latency and Throughput}

Writing audit transactions to a distributed ledger introduces
latency. A triage agent writing an on-chain transaction for each
decision step faces throughput constraints that delay
time-sensitive interventions. Off-chain computation with on-chain
commitment, periodic batch anchoring of audit summaries, and
layer-2 scaling protocols each offer partial solutions, but all
require careful design to preserve the immutability properties
that justify on-chain governance in the first place.

\subsection*{Interoperability with Existing Health Information
Systems}

Healthcare institutions run a heterogeneous mix of electronic
health record systems, laboratory information systems, and
clinical data warehouses built before blockchain-based
architectures existed. Integrating on-chain governance requires
FHIR-compliant data exchange, HL7 interface standards, and
institutional commitment to new auditing protocols. This is
substantial integration engineering, and it cannot be wished
away.

\subsection*{GDPR and the Right to Erasure}

Blockchain immutability conflicts directly with the GDPR right to
erasure. A patient invoking the right to deletion cannot have
data removed from an immutable ledger. One architectural approach
stores only cryptographic commitments on-chain while keeping
personal data off-chain. Erasure is implemented by deleting the
off-chain data and rendering the commitment unresolvable,
preserving both the audit record's integrity and the practical
effect of erasure. Bai et~al.\ demonstrate this hybrid model in
practice for smart healthcare systems, using IPFS for off-chain
storage of personally identifiable and protected health
information while maintaining only encrypted references on the
blockchain ledger, implementing GDPR Article~17 compliance by
enabling deletion from IPFS storage without compromising the
integrity of the on-chain audit record.\cite{bai2022} The broader
legal adequacy of these hybrid approaches across different
regulatory jurisdictions remains an open question.

\subsection*{Institutional Adoption and Governance}

Distributed ledger governance requires consortium formation:
agreement on ledger architecture, consensus mechanisms, access
control, and governance of the governance layer itself. Hospitals
and health systems accustomed to managing their own data
infrastructure resist participating in a shared ledger they do
not control unilaterally. Building institutional trust and
regulatory legitimacy for cross-institutional blockchain-based AI
governance is a sociotechnical problem, and technical
architecture alone does not solve it.

\subsection*{Clinician Trust Calibration}

Governance frameworks typically focus on system-level
accountability. But agentic AI governance has a human layer too.
Clinicians who are over-reliant on AI agent recommendations and
clinicians who distrust them both pose risks. In this context,
calibration refers not only to probability reliability but to the
alignment between a clinician's confidence in an AI
recommendation and the system's actual track record. Frameworks
that improve technical trustworthiness also need to invest in
clinician education about when and how to weight AI output
appropriately.

%% ============================================================
%% SECTION 7 -- RESEARCH AGENDA
%% ============================================================
\section*{Research Agenda}

This framework is a conceptual contribution. Empirical validation
is required. We identify five priorities.

\textbf{Priority 1: Empirical calibration evaluation of deployed
clinical AI.}
The reliability layer requires calibration measurement and
documentation for each agent decision node. Large-scale empirical
studies measuring ECE, Brier scores, and conformal coverage rates
across deployed clinical AI systems, including agentic pipelines,
would establish the current state and identify where governance
intervention is most urgent.

\textbf{Priority 2: Smart contract frameworks for clinical AI
ethical constraints.}
Encoding governance rules in smart contract code requires
healthcare-specific contract templates validated against clinical
governance requirements, especially for escalation triggers,
consent enforcement, and calibration drift alerts. Piloting these
in simulated clinical environments would produce evidence of
operational feasibility.

\textbf{Priority 3: Cross-institutional federated governance
pilots.}
Multi-institutional pilots are needed to evaluate feasibility,
latency, interoperability, and governance dynamics for the
federated learning architecture described above. These pilots
would also generate the diverse external validation datasets
needed to evaluate calibration stability and subgroup fairness
across real institutional variation.

\textbf{Priority 4: Patient-facing uncertainty communication
standards.}
If conformal prediction is to be used to communicate uncertainty
to patients and clinicians in ways that support informed consent
and meaningful oversight, standards are needed for how that
uncertainty is displayed, interpreted, and acted upon in clinical
interfaces. Human factors research bridging the technical and
clinical dimensions of uncertainty communication is a necessary
investment.

\textbf{Priority 5: Prospective evaluation of DLT audit trails
in adverse event investigation.}
The value of blockchain-based audit trails is ultimately
demonstrated by studying their use in actual adverse event
reviews. Prospective studies comparing the completeness,
accuracy, and utility of on-chain versus conventional audit
records in clinical AI incident investigations would provide the
evidence base needed to justify the investment DLT adoption in
healthcare AI governance requires.

%% ============================================================
%% SECTION 8 -- CONCLUSION
%% ============================================================
\section*{Conclusion}

Agentic AI in healthcare is not a future prospect. It is here.
Systems that autonomously monitor, reason, and act on behalf of
patients are in clinical use now, and their capabilities are
expanding. The governance frameworks available to health systems,
regulators, and patients were built for a simpler picture: a
model produces a prediction, a clinician evaluates it, a patient
receives care. AI agents now make sequential decisions at machine
speed. Those frameworks leave real accountability gaps open.

The five-layer trustworthiness governance framework proposed here
addresses those gaps with measurable, testable criteria.
Reliability evaluation, fairness auditing, provenance
verification, consent management, and deployment readiness
scoring each carry specific thresholds rather than aspirational
language. Distributed ledger technologies provide the
accountability infrastructure that makes the framework
enforceable: immutable audit trails, smart contract-enforced
constraints, decentralized consent registries, and on-chain model
provenance records transform governance from written policy into
operational architecture.

Trustworthiness in agentic healthcare AI is not something
certified once at deployment and then assumed. It must be earned
continuously, measured rigorously, and verified independently.
The framework here offers a concrete path toward that goal, and a
research agenda for the empirical work required to get there.

%% ============================================================
%% CLOSING SECTIONS (BHTY required order)
%% ============================================================

\section*{Acknowledgments}
The author acknowledges the growing body of clinical AI
governance scholarship that informed this perspective,
particularly the TRIPOD+AI, DECIDE-AI, and emerging conformal
prediction frameworks for clinical uncertainty quantification.
Portions of this manuscript were drafted with the assistance of
Claude (Anthropic), an AI language model, used for structural
drafting, editorial refinement, and reference formatting. The
author takes full responsibility for all content, interpretations,
and conclusions.

\section*{Financial/Non-Financial Disclosures}
\textbf{Funding:} None.

\textbf{Conflicts of interest:} The author declares no financial
or non-financial conflicts of interest relevant to the content of
this manuscript.

\section*{Contributors}
Michael O. Eniolade contributed to conceptualization, framework
development, formal analysis, writing --- original draft, and
writing --- review and editing.

\section*{Data Availability Statement (DAS), Data Sharing,
Reproducibility, and Data Repositories}
No new data were generated or analyzed in this study. This is a
conceptual perspective paper. All supporting evidence is drawn
from published literature cited in the references.

\section*{Application of AI-Generated Text or Related Technology}
Portions of this manuscript were drafted with the assistance of
Claude (Anthropic), an AI language model, used for structural
drafting, editorial refinement, and reference formatting. The
author takes full responsibility for all content, interpretations,
and conclusions presented in this work.

%% ============================================================
%% REFERENCES  (Vancouver numbered, superscript)
%% ============================================================
\section*{References}

\begin{thebibliography}{19}
\setlength{\itemsep}{4pt}

\bibitem{rajpurkar2022}
Rajpurkar P, Chen E, Banerjee O, Topol EJ. AI in health and
medicine. \textit{Nat Med.} 2022;28(1):31--8.
\url{https://doi.org/10.1038/s41591-021-01614-0}

\bibitem{gorenshtein2025}
Gorenshtein A, Omar M, Glicksberg BS, Nadkarni GN, Klang E.
AI Agents in Clinical Medicine: A Systematic Review.
\textit{medRxiv.} 2025;2025.08.22.25334232.
\url{https://doi.org/10.1101/2025.08.22.25334232}

\bibitem{collins2024tripod}
Collins GS, Moons KGM, Dhiman P, et al. TRIPOD+AI statement:
updated guidance for reporting clinical prediction models that
use regression or machine learning methods. \textit{BMJ.}
2024;385:e078378.
\url{https://doi.org/10.1136/bmj-2023-078378}

\bibitem{vasey2022}
Vasey B, Nagendran M, Campbell B, et al. Reporting guideline
for the early stage clinical evaluation of decision support
systems driven by artificial intelligence: DECIDE-AI.
\textit{BMJ.} 2022;377:e070904.
\url{https://doi.org/10.1136/bmj-2022-070904}

\bibitem{fda2021}
U.S. Food and Drug Administration. Artificial
Intelligence/Machine Learning (AI/ML)-Based Software as a
Medical Device (SaMD) Action Plan. Silver Spring, MD: FDA; 2021.

\bibitem{euaiact2024}
European Parliament and Council. Regulation (EU) 2024/1689 of
the European Parliament and of the Council --- Artificial
Intelligence Act. \textit{Official Journal of the European
Union.} 2024.

\bibitem{guo2017}
Guo C, Pleiss G, Sun Y, Weinberger KQ. On calibration of
modern neural networks. \textit{Proc 34th Int Conf Mach Learn.}
2017;70:1321--30.

\bibitem{angelopoulos2023}
Angelopoulos AN, Bates S. A gentle introduction to conformal
prediction and distribution-free uncertainty quantification.
\textit{Found Trends Mach Learn.} 2023;16(4):494--591.
\url{https://doi.org/10.1561/2200000101}

\bibitem{choudhury2020}
Choudhury A, Asan O. Role of artificial intelligence in patient
safety outcomes: systematic literature review. \textit{JMIR Med
Inform.} 2020;8(7):e18599.
\url{https://doi.org/10.2196/18599}

\bibitem{zuiderwijk2021}
Zuiderwijk A, Chen YC, Salem F. Implications of the use of
artificial intelligence in public governance: A systematic
literature review and a research agenda. \textit{Gov Inf Q.}
2021;38:101577.
\url{https://doi.org/10.1016/j.giq.2021.101577}

\bibitem{haris2024}
Haris M, Raveendra K, Travlos CK, Lewington A, Wu J,
Shuweidhi F, Nadarajah R, Gale CP. Prediction of incident
chronic kidney disease in community-based electronic health
records: a systematic review and meta-analysis.
\textit{Clin Kidney J.} 2024;17(5):sfae098.
\url{https://doi.org/10.1093/ckj/sfae098}

\bibitem{prinzi2025}
Campagner A, Ciucci D, Cabitza F. Modeling the unknown: on the
role of uncertainty in machine learning for healthcare.
\textit{Comput Methods Programs Biomed.} 2025;260:108576.
\url{https://doi.org/10.1016/j.cmpb.2024.108576}

\bibitem{mittelstadt2016}
Mittelstadt B, Allo P, Taddeo M, Wachter S, Floridi L. The
ethics of algorithms: mapping the debate. \textit{Big Data Soc.}
2016;3(2):2053951716679679.
\url{https://doi.org/10.1177/2053951716679679}

\bibitem{collins2024eval}
Collins GS, Dhiman P, Ma J, et al. Evaluation of clinical
prediction models (part 1): from development to external
validation. \textit{BMJ.} 2024;384:e074819.
\url{https://doi.org/10.1136/bmj-2023-074819}

\bibitem{dershem2025}
Dershem M, Riley~III J, Venkataraman M, Nasr J, Hussain AS.
2026 Healthcare Predictions: AI, Blockchain, and the Rise of
Decentralized Innovation. \textit{BHTY.} 2025;8(3).
\url{https://doi.org/10.30953/bhty.v8.475}

\bibitem{ullah2024}
Ullah F, He J, Zhu N, Wajahat A, Nazir A, Qureshi S,
Pathan MS, Dev S. Blockchain-enabled EHR access auditing:
Enhancing healthcare data security. \textit{Heliyon.}
2024;10(16):e34407.
\url{https://doi.org/10.1016/j.heliyon.2024.e34407}

\bibitem{liu2026}
Liu J, Hu X. Blockchain meets AI in healthcare: a review of
convergent technologies for digital health transformation.
\textit{Front Blockchain.} 2026;9:1766092.
\url{https://doi.org/10.3389/fbloc.2026.1766092}

\bibitem{rieke2020}
Rieke N, Hancox J, Li W, et al. The future of digital health
with federated learning. \textit{npj Digit Med.}
2020;3(1):119.
\url{https://doi.org/10.1038/s41746-020-00323-1}

\bibitem{bai2022}
Bai P, Kumar S, Kumar K, Kaiwartya O, Mahmud M, Lloret J.
GDPR Compliant Data Storage and Sharing in Smart Healthcare
System: A Blockchain-Based Solution. \textit{Electronics.}
2022;11(20):3311.
\url{https://doi.org/10.3390/electronics11203311}

\end{thebibliography}

\end{document}
